% ============================================================
% CAPÍTULOS SEM NUMERAÇÃO
% ============================================================

\newcommand{\capitulopreliminar}[1]{%
    \cleardoublepage
    \chapter*{#1}
    \addcontentsline{toc}{chapter}{#1}
    \markboth{#1}{#1}
}

% ==========================================================
% TEOREMAS
% ==========================================================


% ==========================================================
% ARMAZENAMENTO DE SOLUÇÃO
% ==========================================================

\newcounter{contsolucao}

\newenvironment{solucoes}[1]{%
    \setcounter{contsolucao}{0}%
    \def\totalSolucoes{#1}%
    \par\medskip
}{%
    \ifnum\value{contsolucao}>0
        \end{caixasolucao}
    \fi
    \par\medskip
}

% ==================================================
% COMANDO DE SOLUÇÃO
% ==================================================

\newcommand{\solucao}{%
    \ifnum\value{contsolucao}>0
        \end{caixasolucao}
    \fi
    \stepcounter{contsolucao}%
    \begin{caixasolucao}
    \textbf{\textcolor{VermelhoSolucao}{Solução%
        \ifnum\totalSolucoes>1
            \space
            \ifnum\value{contsolucao}<10
                0%
            \fi
            \arabic{contsolucao}%
        \fi
    }}%
    \par\smallskip
}
% ==================================================
% CONTADOR DE QUESTÕES
% ==================================================

\newcounter{numeroquestao}

\renewcommand{\thenumeroquestao}{%
    \ifnum\value{numeroquestao}<10
        0\arabic{numeroquestao}%
    \else
        \arabic{numeroquestao}%
    \fi
}

\makeatletter

\let\sectionoriginal\section

\renewcommand{\section}{%
    \@ifstar
        {\sectionsemnumero}
        {\sectioncomnumero}%
}

\newcommand{\sectioncomnumero}[1]{%
    \setcounter{numeroquestao}{0}%
    \sectionoriginal{#1}%
}

\newcommand{\sectionsemnumero}[1]{%
    \setcounter{numeroquestao}{0}%
    \sectionoriginal*{#1}%
}

\makeatother

% ============================================================
% ESTILO DAS QUESTÕES
%=============================================================

\renewcommand{\questao}[1]{%
    \refstepcounter{numeroquestao}
    \noindent
    \hangindent=0cm
    \hangafter=-1
    \makebox[2.85cm][l]{
    \begin{tikzpicture}[baseline=-.25cm]
        \node[
            fill=AzulPrincipal,
            text=white,
            rounded corners=3pt,
            minimum width=1cm,
            minimum height=0.8cm,
            font=\bfseries
        ] (numero) {Questão \thenumeroquestao};
        \draw[
            Turquesa,
            line width=2pt
        ]
        ([xshift=.3cm, yshift=.37cm]numero.east) -- ++(\dimexpr\textwidth-2.95cm\relax,0);
        
    \end{tikzpicture}%
    }
    #1
    \par 
    \vspace{.3cm}
    }


% =====================================================================
% PÁGINA DEDICADA AOS AUTORES
% =====================================================================

\newcommand{\paginaautores}{%
    \cleardoublepage
    \thispagestyle{empty}

    \begin{center}

        \vspace*{1cm}

        {\Large\bfseries\color{AzulPrincipal}
        AUTORES}

        \vspace{0.3cm}

        {\color{Turquesa}\rule{5cm}{2pt}}

        \vspace{1.2cm}

        % Autor 1
        {\large\bfseries\color{AzulPrincipal}
        Jadiel José Araújo da Silva Júnior}

        \vspace{0.2cm}

        {\normalsize
        Licenciado em Matemática - UFAL\\
        Mestre em Matemática - UFAL}

        \vspace{1cm}

        % Autor 2
        {\large\bfseries\color{AzulPrincipal}
        Samuel de Sousa Apolinário}

        \vspace{0.2cm}

        {\normalsize
        Licenciado em Matemática - UNEAL\\
        Especialista em Estatística Aplicada - FFOCUS\\
        Especialista em Matemática Orientada a Computação Gráfica - FFOCUS}

        \vspace{1cm}

        % Autor 3
        {\large\bfseries\color{AzulPrincipal}
        Gabriel Dória de Mendonça}

        \vspace{0.2cm}

        {\normalsize
        Licenciado em Matemática - UFAL}

        \vfill

        {\color{Turquesa}\rule{5cm}{1pt}}

        \vspace{0.3cm}

        {\small
        Banco de Questões de Matemática}

    \end{center}

    \cleardoublepage
}